\documentclass[UTF8,12pt]{ctexart}
\usepackage[a4paper,margin=24mm]{geometry}
\usepackage{amsmath,amssymb,bm,graphicx,adjustbox}
\newcommand{\fitmath}[1]{\begin{adjustbox}{max width=\linewidth}$\displaystyle #1$\end{adjustbox}}
\setlength{\parindent}{0pt}
\setlength{\parskip}{.7em}
\sloppy
\allowdisplaybreaks
\begin{document}
\section*{2014 年考研数学三 · 真题}
\subsection*{原 PDF 第 26 页}

\subsection*{2014 年全国硕士研究生招生考试试题}

\subsection*{一、选择题（本题共 8 小题，每小题 4 分，共 32 分。在每小题给出的四个选项中，只有一项符合题目要求，把所选项前的字母填在题后的括号内。）}

(1) \(\displaystyle \lim_{n\to\infty}a_n=a\)，且 \(\displaystyle a\ne0\)，则当 \(\displaystyle n\) 充分大时有（　）。


(A) \(\displaystyle |a_n|>\dfrac{\displaystyle |a|}{\displaystyle 2}\)。 (B) \(\displaystyle |a_n|<\dfrac{\displaystyle |a|}{\displaystyle 2}\)。 (C) \(\displaystyle a_n>a-\dfrac{\displaystyle 1}{\displaystyle n}\)。 (D) \(\displaystyle a_n<a+\dfrac{\displaystyle 1}{\displaystyle n}\)。


(2) 下列曲线中有渐近线的是（　）。


(A) \(\displaystyle y=x+\sin x\)。 (B) \(\displaystyle y=x^2+\sin x\)。 (C) \(\displaystyle y=x+\sin\dfrac{\displaystyle 1}{\displaystyle x}\)。 (D) \(\displaystyle y=x^2+\sin\dfrac{\displaystyle 1}{\displaystyle x}\)。


(3) 设 \(\displaystyle p(x)=a+bx+cx^2+dx^3\)。当 \(\displaystyle x\to0\) 时，若 \(\displaystyle p(x)-\tan x\) 是比 \(\displaystyle x^3\) 高阶的无穷小量，则下列选项中错误的是（　）。


(A) \(\displaystyle a=0\)。 (B) \(\displaystyle b=1\)。 (C) \(\displaystyle c=0\)。 (D) \(\displaystyle d=\dfrac{\displaystyle 1}{\displaystyle 6}\)。


(4) 设函数 \(\displaystyle f(x)\) 具有 2 阶导数，\(\displaystyle g(x)=f(0)(1-x)+f(1)x\)，则在区间 \(\displaystyle [0,\allowbreak 1]\) 上（　）。


(A) 当 \(\displaystyle f^{\prime}(x)\geq0\) 时，\(\displaystyle f(x)\geq g(x)\)。 (B) 当 \(\displaystyle f^{\prime}(x)\geq0\) 时，\(\displaystyle f(x)\leq g(x)\)。 (C) 当 \(\displaystyle f^{\prime\prime}(x)\geq0\) 时，\(\displaystyle f(x)\geq g(x)\)。 (D) 当 \(\displaystyle f^{\prime\prime}(x)\geq0\) 时，\(\displaystyle f(x)\leq g(x)\)。


(5) 行列式 \(\displaystyle \begin{vmatrix}0&a&b&0\\a&0&0&b\\0&c&d&0\\c&0&0&d\end{vmatrix}=\)（　）。


(A) \(\displaystyle (ad-bc)^2\)。 (B) \(\displaystyle -(ad-bc)^2\)。 (C) \(\displaystyle a^2d^2-b^2c^2\)。 (D) \(\displaystyle b^2c^2-a^2d^2\)。


(6) 设 \(\displaystyle \alpha_1,\allowbreak \alpha_2,\allowbreak \alpha_3\) 均为 3 维向量，则对任意常数 \(\displaystyle k,\allowbreak l\)，向量组 \(\displaystyle \alpha_1+k\alpha_3,\allowbreak \alpha_2+l\alpha_3\) 线性无关是向量组 \(\displaystyle \alpha_1,\allowbreak \alpha_2,\allowbreak \alpha_3\) 线性无关的（　）。


(A) 必要非充分条件。 (B) 充分非必要条件。 (C) 充分必要条件。 (D) 既非充分也非必要条件。


(7) 设随机事件 \(\displaystyle A\) 与 \(\displaystyle B\) 相互独立，且 \(\displaystyle P(B)=0.5,\allowbreak P(A-B)=0.3\)，则 \(\displaystyle P(B-A)=\)（　）。


(A) 0.1。 (B) 0.2。 (C) 0.3。 (D) 0.4。


(8) 设 \(\displaystyle X_1,\allowbreak X_2,\allowbreak X_3\) 为来自正态总体 \(\displaystyle N(0,\allowbreak \sigma^2)\) 的简单随机样本，则统计量 \(\displaystyle S=\dfrac{\displaystyle X_1-X_2}{\displaystyle \sqrt2|X_3|}\) 服从的分布为（　）。


(A) \(\displaystyle F(1,\allowbreak 1)\)。 (B) \(\displaystyle F(2,\allowbreak 1)\)。 (C) \(\displaystyle t(1)\)。 (D) \(\displaystyle t(2)\)。


\subsection*{二、填空题（本题共 6 小题，每小题 4 分，共 24 分，把答案填在题中横线上。）}

(9) 设某商品的需求函数为 \(\displaystyle Q=40-2P\)（\(\displaystyle P\) 为商品的价格），则该商品的边际收益为\_\_\_\_\_\_。


(10) 设 \(\displaystyle D\) 是由曲线 \(\displaystyle xy+1=0\) 与直线 \(\displaystyle y+x=0\) 及 \(\displaystyle y=2\) 围成的有界区域，则 \(\displaystyle D\) 的面积为\_\_\_\_\_\_。


\clearpage

\subsection*{原 PDF 第 27 页}

(11) 设 \(\displaystyle \int_0^a xe^{2x}\,dx=\dfrac{\displaystyle 1}{\displaystyle 4}\)，则 \(\displaystyle a=\)\_\_\_\_\_\_。


(12) 二次积分 \(\displaystyle \int_0^1dy\displaystyle \int_y^1\left(\dfrac{\displaystyle e^{x^2}}{\displaystyle x}-e^{y^2}\right)dx=\)\_\_\_\_\_\_。


(13) 设二次型 \(\displaystyle f(x_1,\allowbreak x_2,\allowbreak x_3)=x_1^2-x_2^2+2ax_1x_3+4x_2x_3\) 的负惯性指数为 1，则 \(\displaystyle a\) 的取值范围是\_\_\_\_\_\_。


(14) 设总体 \(\displaystyle X\) 的概率密度为 \(\displaystyle f(x;\theta)=\begin{cases}\dfrac{\displaystyle 2x}{\displaystyle 3\theta^2},\allowbreak &\theta<x<2\theta,\allowbreak \\0,\allowbreak &\text{其他},\allowbreak \end{cases}\) 其中 \(\displaystyle \theta\) 是未知参数，\(\displaystyle X_1,\allowbreak X_2,\allowbreak \ldots,\allowbreak X_n\) 为来自总体 \(\displaystyle X\) 的简单随机样本。若 \(\displaystyle E\left(c\displaystyle \sum_{i=1}^nX_i^2\right)=\theta^2\)，则 \(\displaystyle c=\)\_\_\_\_\_\_。


\subsection*{三、解答题（本题共 9 小题，共 94 分，解答应写出文字说明、证明过程或演算步骤。）}

(15)（本题满分 10 分）求极限 \(\displaystyle \lim_{x\to+\infty}\dfrac{\displaystyle \int_1^x[t^2(e^{1/t}-1)-t]\,dt}{\displaystyle x^2\ln(1+1/x)}\)。


(16)（本题满分 10 分）设平面区域 \(\displaystyle D=\{(x,\allowbreak y)\mid1\leq x^2+y^2\leq4,\allowbreak \ x\geq0,\allowbreak \ y\geq0\}\)，计算 \(\displaystyle \iint_D\dfrac{\displaystyle x\sin(\pi\sqrt{x^2+y^2})}{\displaystyle x+y}\,dx\,dy\)。


\clearpage

\subsection*{原 PDF 第 28 页}

(17)（本题满分 10 分）设函数 \(\displaystyle f(u)\) 具有连续导数，且 \(\displaystyle z=f(e^x\cos y)\) 满足 \(\displaystyle \cos y\dfrac{\displaystyle \partial z}{\displaystyle \partial x}-\sin y\dfrac{\displaystyle \partial z}{\displaystyle \partial y}=(4z+e^x\cos y)e^x\)。若 \(\displaystyle f(0)=0\)，求 \(\displaystyle f(u)\) 的表达式。


(18)（本题满分 10 分）求幂级数 \(\displaystyle \sum_{n=0}^{\infty}(n+1)(n+3)x^n\) 的收敛域及和函数。


(19)（本题满分 10 分）设函数 \(\displaystyle f(x),\allowbreak g(x)\) 在区间 \(\displaystyle [a,\allowbreak b]\) 上连续，且 \(\displaystyle f(x)\) 单调增加，\(\displaystyle 0\leq g(x)\leq1\)。证明：


（I）\(\displaystyle 0\leq\displaystyle \int_a^x g(t)\,dt\leq x-a,\allowbreak \ x\in[a,\allowbreak b]\)；（II）\(\displaystyle \int_a^{a+\displaystyle \int_a^b g(t)\,dt} f(x)\,dx\leq\displaystyle \int_a^b f(x)g(x)\,dx\)。


(20)（本题满分 11 分）设矩阵 \(\displaystyle A=\begin{pmatrix}1&-2&3&-4\\0&1&-1&1\\1&2&0&-3\end{pmatrix}\)，\(\displaystyle E\) 为 3 阶单位矩阵。


（I）求方程组 \(\displaystyle Ax=0\) 的一个基础解系；（II）求满足 \(\displaystyle AB=E\) 的所有矩阵 \(\displaystyle B\)。


\clearpage

\subsection*{原 PDF 第 29 页}

(21)（本题满分 11 分）证明 \(\displaystyle n\) 阶矩阵 \(\displaystyle \begin{pmatrix}1&1&\cdots&1\\1&1&\cdots&1\\\vdots&\vdots&&\vdots\\1&1&\cdots&1\end{pmatrix}\) 与 \(\displaystyle \begin{pmatrix}0&\cdots&0&1\\0&\cdots&0&2\\\vdots&&\vdots&\vdots\\0&\cdots&0&n\end{pmatrix}\) 相似。


(22)（本题满分 11 分）设随机变量 \(\displaystyle X\) 的概率分布为 \(\displaystyle P\{X=1\}=P\{X=2\}=\dfrac{\displaystyle 1}{\displaystyle 2}\)。在给定 \(\displaystyle X=i\) 的条件下，随机变量 \(\displaystyle Y\) 服从均匀分布 \(\displaystyle U(0,\allowbreak i)\)（\(\displaystyle i=1,\allowbreak 2\)）。


（I）求 \(\displaystyle Y\) 的分布函数 \(\displaystyle F_Y(y)\)；（II）求 \(\displaystyle E(Y)\)。


(23)（本题满分 11 分）设随机变量 \(\displaystyle X,\allowbreak Y\) 的概率分布相同，\(\displaystyle X\) 的概率分布为 \(\displaystyle P\{X=0\}=\dfrac{\displaystyle 1}{\displaystyle 3},\allowbreak \ P\{X=1\}=\dfrac{\displaystyle 2}{\displaystyle 3}\)，且 \(\displaystyle X\) 与 \(\displaystyle Y\) 的相关系数 \(\displaystyle \rho_{XY}=\dfrac{\displaystyle 1}{\displaystyle 2}\)。


（I）求 \(\displaystyle (X,\allowbreak Y)\) 的概率分布；（II）求 \(\displaystyle P\{X+Y\leq1\}\)。


\clearpage
\end{document}
